\chapter{Winding Numbers and the Argument Principle}
\label{ch:iv12}

Winding numbers measure how closed curves wrap around points, while the argument principle converts winding into an exact count of zeros and poles. This chapter is the bridge from contour integration to global counting theorems.

\section*{Learning goals}
After this chapter, the reader should be able to:
\begin{itemize}
\item compute winding numbers of explicit closed curves around specified points;
\item apply the argument principle to count zeros minus poles inside a contour;
\item derive zero/pole counts from the logarithmic derivative \(f'/f\);
\item track multiplicities correctly in argument-principle calculations;
\item relate change of argument along a contour to topological winding;
\item identify failures when the contour passes through zeros or poles;
\end{itemize}
\section*{Conceptual roadmap}
\[
\boxed{\text{local holomorphic structure}\;\longrightarrow\;\text{integral or series representation}\;\longrightarrow\;\text{global consequence}.}
\]

\section{Index of a closed curve}
For a closed curve $\Gamma$ avoiding $a$, the index $\operatorname{Ind}(\Gamma,a)=\frac{1}{2\pi i}\int_\Gamma \frac{dz}{z-a}$ records the net winding around $a$.

% BEGIN VOL04-EXPANSION IV12-example-01
\begin{example}[Winding number of a translated circle]\label{ex:iv12-expand-01}
Let \(\gamma(t)=2+3e^{it}\), \(0\le t\le2\pi\). Around the point \(a=2\),
\[
\operatorname{Ind}(\gamma,2)=\frac{1}{2\pi i}\int_\gamma\frac{dz}{z-2}=1.
\]
Around \(a=6\), the point lies outside the circle and the index is zero. The index records how the curve winds around a chosen point, not merely the shape of the curve itself.
\end{example}
% END VOL04-EXPANSION IV12-example-01
\section{Homotopy invariance}
The winding number is unchanged under deformations of the curve that avoid the reference point.

\section{Local constancy}
As a function of $a$ off the curve, the winding number is integer-valued and locally constant on each component of the complement.

\section{Logarithmic derivative}
For meromorphic $f$, the quotient $f'/f$ has residues equal to zero multiplicities and negatives of pole multiplicities.

% BEGIN VOL04-EXPANSION IV12-example-02
\begin{example}[Argument principle for a polynomial]\label{ex:iv12-expand-02}
Take \(f(z)=z^3-1\) on \(|z|=2\). There are no zeros or poles on the contour, and all three roots lie inside. Therefore
\[
\frac{1}{2\pi i}\int_{|z|=2}\frac{f'(z)}{f(z)}\,dz=3.
\]
Equivalently, the image curve \(f(2e^{it})\) winds three times around the origin.
\end{example}
% END VOL04-EXPANSION IV12-example-02
\section{Argument principle}
The contour integral of $f'/f$ counts zeros minus poles inside a contour, weighted by multiplicity.

\section{Change of argument}
The argument principle can be interpreted as total variation of the argument of $f(\Gamma)$.

\section{Counting polynomial roots}
Contours can count roots in regions without locating them explicitly.

% BEGIN VOL04-EXPANSION IV12-example-03
\begin{example}[Zeros minus poles]\label{ex:iv12-expand-03}
For \(f(z)=(z-1)^2/(z+1)\) on \(|z|=2\), the contour encloses a double zero at \(1\) and a simple pole at \(-1\). The argument principle gives
\[
\frac{1}{2\pi i}\int_{|z|=2}\frac{f'}{f}\,dz=2-1=1.
\]
Multiplicity and pole order enter with opposite signs.
\end{example}
% END VOL04-EXPANSION IV12-example-03
\section{General residue formulation}
The winding-number form of the residue theorem allows multiply wound contours and non-simple geometry.

\section{Boundary hypotheses}
The argument principle requires the contour to avoid zeros and poles of the meromorphic function.

\section{Topological synthesis}
Winding, residues, multiplicity, and image curves are different faces of the same integer-valued invariant.

\section{Core structural results}

\begin{theorem}[Integer-valued index]\label{thm:iv12-01}
For a closed piecewise $C^1$ curve $\Gamma$ and $a\notin\Gamma$, $\operatorname{Ind}(\Gamma,a)$ is an integer.
\begin{proof}
Locally choose a branch of the logarithm along successive short pieces of the curve. The integral of $(z-a)^{-1}$ is the net change of logarithm, and after one circuit the logarithm can change only by an integral multiple of $2\pi i$.
\end{proof}
\end{theorem}

\begin{theorem}[Argument principle]\label{thm:iv12-02}
If $f$ is meromorphic inside and on $\Gamma$, with no zeros or poles on $\Gamma$, then $\frac{1}{2\pi i}\int_\Gamma f'(z)/f(z)\,dz=N-P$, counting zeros and poles with multiplicity in the simple winding-one case.
\begin{proof}
The logarithmic derivative has a simple pole at every zero with residue equal to its multiplicity and at every pole with residue minus its order. Apply the residue theorem.
\end{proof}
\end{theorem}

\begin{theorem}[Homotopy invariance of index]\label{thm:iv12-03}
If two closed curves are homotopic in $\mathbb C\setminus\{a\}$, then they have the same winding number about $a$.
\begin{proof}
Apply homotopy invariance of contour integrals to the holomorphic function $(z-a)^{-1}$ on the punctured plane.
\end{proof}
\end{theorem}

\begin{theorem}[Local constancy of index]\label{thm:iv12-04}
The map $a\mapsto\operatorname{Ind}(\Gamma,a)$ is constant on every connected component of $\mathbb C\setminus\Gamma$.
\begin{proof}
The integral depends continuously on $a$ off the curve, while its values are integers. A continuous integer-valued function is locally constant, hence constant on connected components.
\end{proof}
\end{theorem}

\section{Worked examples}

\begin{example}[Index of a closed curve diagnostic]\label{ex:iv12-worked-01}
Give a rigorous worked analysis of Index of a closed curve. State the relevant holomorphic, contour, series, or singularity hypothesis and derive the central conclusion. For a closed curve $\Gamma$ avoiding $a$, the index $\operatorname{Ind}(\Gamma,a)=\frac{1}{2\pi i}\int_\Gamma \frac{dz}{z-a}$ records the net winding around $a$. The decisive step is to use the complex-analytic structure globally rather than reason from real-variable intuition alone.
\end{example}

\begin{example}[Homotopy invariance diagnostic]\label{ex:iv12-worked-02}
Give a rigorous worked analysis of Homotopy invariance. State the relevant holomorphic, contour, series, or singularity hypothesis and derive the central conclusion. The winding number is unchanged under deformations of the curve that avoid the reference point. The decisive step is to use the complex-analytic structure globally rather than reason from real-variable intuition alone.
\end{example}

\begin{example}[Local constancy diagnostic]\label{ex:iv12-worked-03}
Give a rigorous worked analysis of Local constancy. State the relevant holomorphic, contour, series, or singularity hypothesis and derive the central conclusion. As a function of $a$ off the curve, the winding number is integer-valued and locally constant on each component of the complement. The decisive step is to use the complex-analytic structure globally rather than reason from real-variable intuition alone.
\end{example}

\begin{example}[Logarithmic derivative diagnostic]\label{ex:iv12-worked-04}
Give a rigorous worked analysis of Logarithmic derivative. State the relevant holomorphic, contour, series, or singularity hypothesis and derive the central conclusion. For meromorphic $f$, the quotient $f'/f$ has residues equal to zero multiplicities and negatives of pole multiplicities. The decisive step is to use the complex-analytic structure globally rather than reason from real-variable intuition alone.
\end{example}

\section{Solved dossiers}
The dossiers are canonical solved problems. The provenance ledger separately records which retained corpus rules guided each topic and which dossiers were newly designed.

\begin{problem}[Index of a closed curve diagnostic]\label{prob:iv12-dossier-01}
Give a rigorous worked analysis of Index of a closed curve. State the relevant holomorphic, contour, series, or singularity hypothesis and derive the central conclusion.
\end{problem}
\begin{solution}
For a closed curve $\Gamma$ avoiding $a$, the index $\operatorname{Ind}(\Gamma,a)=\frac{1}{2\pi i}\int_\Gamma \frac{dz}{z-a}$ records the net winding around $a$. The decisive step is to use the complex-analytic structure globally rather than reason from real-variable intuition alone.
\end{solution}

\begin{problem}[Homotopy invariance diagnostic]\label{prob:iv12-dossier-02}
Give a rigorous worked analysis of Homotopy invariance. State the relevant holomorphic, contour, series, or singularity hypothesis and derive the central conclusion.
\end{problem}
\begin{solution}
The winding number is unchanged under deformations of the curve that avoid the reference point. The decisive step is to use the complex-analytic structure globally rather than reason from real-variable intuition alone.
\end{solution}

\begin{problem}[Local constancy diagnostic]\label{prob:iv12-dossier-03}
Give a rigorous worked analysis of Local constancy. State the relevant holomorphic, contour, series, or singularity hypothesis and derive the central conclusion.
\end{problem}
\begin{solution}
As a function of $a$ off the curve, the winding number is integer-valued and locally constant on each component of the complement. The decisive step is to use the complex-analytic structure globally rather than reason from real-variable intuition alone.
\end{solution}

\begin{problem}[Logarithmic derivative diagnostic]\label{prob:iv12-dossier-04}
Give a rigorous worked analysis of Logarithmic derivative. State the relevant holomorphic, contour, series, or singularity hypothesis and derive the central conclusion.
\end{problem}
\begin{solution}
For meromorphic $f$, the quotient $f'/f$ has residues equal to zero multiplicities and negatives of pole multiplicities. The decisive step is to use the complex-analytic structure globally rather than reason from real-variable intuition alone.
\end{solution}

\begin{problem}[Argument principle diagnostic]\label{prob:iv12-dossier-05}
Give a rigorous worked analysis of Argument principle. State the relevant holomorphic, contour, series, or singularity hypothesis and derive the central conclusion.
\end{problem}
\begin{solution}
The contour integral of $f'/f$ counts zeros minus poles inside a contour, weighted by multiplicity. The decisive step is to use the complex-analytic structure globally rather than reason from real-variable intuition alone.
\end{solution}

\begin{problem}[Change of argument diagnostic]\label{prob:iv12-dossier-06}
Give a rigorous worked analysis of Change of argument. State the relevant holomorphic, contour, series, or singularity hypothesis and derive the central conclusion.
\end{problem}
\begin{solution}
The argument principle can be interpreted as total variation of the argument of $f(\Gamma)$. The decisive step is to use the complex-analytic structure globally rather than reason from real-variable intuition alone.
\end{solution}

\begin{problem}[Counting polynomial roots diagnostic]\label{prob:iv12-dossier-07}
Give a rigorous worked analysis of Counting polynomial roots. State the relevant holomorphic, contour, series, or singularity hypothesis and derive the central conclusion.
\end{problem}
\begin{solution}
Contours can count roots in regions without locating them explicitly. The decisive step is to use the complex-analytic structure globally rather than reason from real-variable intuition alone.
\end{solution}

\begin{problem}[Integer-valued index proof dossier]\label{prob:iv12-dossier-08}
Prove the following theorem-level statement and identify the essential hypothesis: For a closed piecewise $C^1$ curve $\Gamma$ and $a\notin\Gamma$, $\operatorname{Ind}(\Gamma,a)$ is an integer.
\end{problem}
\begin{solution}
Locally choose a branch of the logarithm along successive short pieces of the curve. The integral of $(z-a)^{-1}$ is the net change of logarithm, and after one circuit the logarithm can change only by an integral multiple of $2\pi i$. This identifies the mechanism that makes the complex-analytic conclusion stronger than its real-variable analogue.
\end{solution}

\begin{problem}[Argument principle proof dossier]\label{prob:iv12-dossier-09}
Prove the following theorem-level statement and identify the essential hypothesis: If $f$ is meromorphic inside and on $\Gamma$, with no zeros or poles on $\Gamma$, then $\frac{1}{2\pi i}\int_\Gamma f'(z)/f(z)\,dz=N-P$, counting zeros and poles with multiplicity in the simple winding-one case.
\end{problem}
\begin{solution}
The logarithmic derivative has a simple pole at every zero with residue equal to its multiplicity and at every pole with residue minus its order. Apply the residue theorem. This identifies the mechanism that makes the complex-analytic conclusion stronger than its real-variable analogue.
\end{solution}

\begin{problem}[Homotopy invariance of index proof dossier]\label{prob:iv12-dossier-10}
Prove the following theorem-level statement and identify the essential hypothesis: If two closed curves are homotopic in $\mathbb C\setminus\{a\}$, then they have the same winding number about $a$.
\end{problem}
\begin{solution}
Apply homotopy invariance of contour integrals to the holomorphic function $(z-a)^{-1}$ on the punctured plane. This identifies the mechanism that makes the complex-analytic conclusion stronger than its real-variable analogue.
\end{solution}

\begin{problem}[Local constancy of index proof dossier]\label{prob:iv12-dossier-11}
Prove the following theorem-level statement and identify the essential hypothesis: The map $a\mapsto\operatorname{Ind}(\Gamma,a)$ is constant on every connected component of $\mathbb C\setminus\Gamma$.
\end{problem}
\begin{solution}
The integral depends continuously on $a$ off the curve, while its values are integers. A continuous integer-valued function is locally constant, hence constant on connected components. This identifies the mechanism that makes the complex-analytic conclusion stronger than its real-variable analogue.
\end{solution}

\begin{problem}[Chapter synthesis]\label{prob:iv12-dossier-12}
Connect the definitions and principal theorems of this chapter into one reusable complex-analysis strategy.
\end{problem}
\begin{solution}
Winding numbers measure how closed curves wrap around points, while the argument principle converts winding into an exact count of zeros and poles. This chapter is the bridge from contour integration to global counting theorems. The reusable strategy is to identify the analytic domain, choose the correct local expansion or contour representation, apply the structural theorem, and then audit singularities and topology.
\end{solution}

\section{Exercises with complete solutions}

\begin{exercise}\label{exr:iv12-01}
State and apply the central principle behind Index of a closed curve. Give one concrete consequence.
\end{exercise}
\begin{hint}
Compute winding number as \((2\pi i)^{-1}\int_\Gamma dz/(z-a)\) or from net change of argument.
\end{hint}
\begin{solution}
For a closed curve $\Gamma$ avoiding $a$, the index $\operatorname{Ind}(\Gamma,a)=\frac{1}{2\pi i}\int_\Gamma \frac{dz}{z-a}$ records the net winding around $a$. The same principle can then be reused in later contour, residue, conformal, or Riemann-surface arguments.
\end{solution}

\begin{exercise}\label{exr:iv12-02}
State and apply the central principle behind Homotopy invariance. Give one concrete consequence.
\end{exercise}
\begin{hint}
Deform the curve without crossing \(a\); winding number stays unchanged under such homotopies.
\end{hint}
\begin{solution}
The winding number is unchanged under deformations of the curve that avoid the reference point. The same principle can then be reused in later contour, residue, conformal, or Riemann-surface arguments.
\end{solution}

\begin{exercise}\label{exr:iv12-03}
State and apply the central principle behind Local constancy. Give one concrete consequence.
\end{exercise}
\begin{hint}
For the argument principle, apply the residue theorem to \(f'/f\) and interpret residues as multiplicities.
\end{hint}
\begin{solution}
As a function of $a$ off the curve, the winding number is integer-valued and locally constant on each component of the complement. The same principle can then be reused in later contour, residue, conformal, or Riemann-surface arguments.
\end{solution}

\begin{exercise}\label{exr:iv12-04}
State and apply the central principle behind Logarithmic derivative. Give one concrete consequence.
\end{exercise}
\begin{hint}
Count zeros and poles with multiplicity; a double zero contributes two, not one.
\end{hint}
\begin{solution}
For meromorphic $f$, the quotient $f'/f$ has residues equal to zero multiplicities and negatives of pole multiplicities. The same principle can then be reused in later contour, residue, conformal, or Riemann-surface arguments.
\end{solution}

\begin{exercise}\label{exr:iv12-05}
State and apply the central principle behind Argument principle. Give one concrete consequence.
\end{exercise}
\begin{hint}
Make sure \(f\) has no zeros or poles on the contour before applying the principle.
\end{hint}
\begin{solution}
The contour integral of $f'/f$ counts zeros minus poles inside a contour, weighted by multiplicity. The same principle can then be reused in later contour, residue, conformal, or Riemann-surface arguments.
\end{solution}

\begin{exercise}\label{exr:iv12-06}
State and apply the central principle behind Change of argument. Give one concrete consequence.
\end{exercise}
\begin{hint}
Track the change of \(\arg f(\Gamma(t))\); its total increment is \(2\pi(N-P)\).
\end{hint}
\begin{solution}
The argument principle can be interpreted as total variation of the argument of $f(\Gamma)$. The same principle can then be reused in later contour, residue, conformal, or Riemann-surface arguments.
\end{solution}

\begin{exercise}\label{exr:iv12-07}
State and apply the central principle behind Counting polynomial roots. Give one concrete consequence.
\end{exercise}
\begin{hint}
Use winding numbers around individual poles when the contour is not a simple Jordan curve.
\end{hint}
\begin{solution}
Contours can count roots in regions without locating them explicitly. The same principle can then be reused in later contour, residue, conformal, or Riemann-surface arguments.
\end{solution}

\begin{exercise}\label{exr:iv12-08}
State and apply the central principle behind General residue formulation. Give one concrete consequence.
\end{exercise}
\begin{hint}
For polynomial root counts, choose a contour on which the leading behavior makes \(f'/f\) or Rouché comparison manageable.
\end{hint}
\begin{solution}
The winding-number form of the residue theorem allows multiply wound contours and non-simple geometry. The same principle can then be reused in later contour, residue, conformal, or Riemann-surface arguments.
\end{solution}

\section*{Chapter summary}
The chapter now forms part of the canonical complex-analysis chain from holomorphicity through residues, global continuation, and Riemann surfaces.
% BEGIN VOL04-EXPANSION IV12-exercises-01
\section{Graded supplementary exercises}
These exercises extend the protected chapter with a balanced set of computations, proofs, hypothesis tests, applications, and challenges.

\subsection*{Standard computations}

\begin{exercise}[Unit circle index]\label{exr:iv12-expand-01}
Compute the winding number of \(e^{it}\), \(0\le t\le2\pi\), about zero.
\end{exercise}
\begin{hint}
Use the logarithmic integral definition.
\end{hint}
\begin{solution}
The index is one.
\end{solution}

\begin{exercise}[Double traversal]\label{exr:iv12-expand-02}
Compute the winding number of \(e^{2it}\) about zero for \(0\le t\le2\pi\).
\end{exercise}
\begin{hint}
Track the change of argument.
\end{hint}
\begin{solution}
The argument increases by \(4\pi\), so the index is two.
\end{solution}

\begin{exercise}[Clockwise circle]\label{exr:iv12-expand-03}
Find the index of \(e^{-it}\) about zero.
\end{exercise}
\begin{hint}
Orientation matters.
\end{hint}
\begin{solution}
The index is minus one.
\end{solution}

\begin{exercise}[Zero count]\label{exr:iv12-expand-04}
Use the argument principle to count zeros of \(z^4+1\) in \(|z|=2\).
\end{exercise}
\begin{hint}
All four roots have modulus one.
\end{hint}
\begin{solution}
There are four zeros, counted with multiplicity, and no poles.
\end{solution}

\begin{exercise}[Zero minus pole]\label{exr:iv12-expand-05}
For \(f(z)=z^3/(z-1)^2\), compute the argument-principle count on \(|z|=2\).
\end{exercise}
\begin{hint}
Count orders inside.
\end{hint}
\begin{solution}
The result is \(3-2=1\).
\end{solution}

\subsection*{Proofs}

\begin{exercise}[Index is integer]\label{exr:iv12-expand-06}
Explain why the winding number of a closed curve about a point not on the curve is an integer.
\end{exercise}
\begin{hint}
Use a continuous change of argument along the curve.
\end{hint}
\begin{solution}
Choose a continuous argument along the parameter interval. Closedness forces the final argument to differ from the initial one by an integer multiple of \(2\pi\), and that integer is the index.
\end{solution}

\begin{exercise}[Homotopy invariance]\label{exr:iv12-expand-07}
Prove that winding number is unchanged under a deformation avoiding the reference point.
\end{exercise}
\begin{hint}
The index varies continuously with the deformation parameter but takes integer values.
\end{hint}
\begin{solution}
The integral defining the index depends continuously on the homotopy. A continuous integer valued function is locally constant, hence constant on the connected parameter interval.
\end{solution}

\begin{exercise}[Argument principle]\label{exr:iv12-expand-08}
Prove the argument principle from residues of \(f'/f\).
\end{exercise}
\begin{hint}
Find the local residue at a zero or pole.
\end{hint}
\begin{solution}
A zero of order \(m\) gives residue \(m\), while a pole of order \(n\) gives residue \(-n\). The residue theorem sums these local integers to give zeros minus poles.
\end{solution}

\begin{exercise}[Change of argument form]\label{exr:iv12-expand-09}
Relate \(\int f'/f\) to the net change of argument of \(f\) along a contour.
\end{exercise}
\begin{hint}
Parametrize the contour and differentiate a local logarithm.
\end{hint}
\begin{solution}
Along intervals where a logarithm is chosen, \(d\log f=f'/f\,dz\). Its real part returns to the initial value on a closed contour, while the imaginary part changes by the total argument increment.
\end{solution}

\subsection*{Counterexamples and hypothesis tests}

\begin{exercise}[Point on curve]\label{exr:iv12-expand-10}
Can the winding number about \(a\) be defined by the usual integral if \(a\) lies on the curve?
\end{exercise}
\begin{hint}
The kernel has a pole on the path.
\end{hint}
\begin{solution}
No. The integral is singular and the standard index is undefined.
\end{solution}

\begin{exercise}[Zero on contour]\label{exr:iv12-expand-11}
Can the argument principle be applied if \(f\) has a zero on the contour?
\end{exercise}
\begin{hint}
Then \(f'/f\) is singular on the path.
\end{hint}
\begin{solution}
No. Zeros and poles must be absent from the contour.
\end{solution}

\begin{exercise}[Ignoring multiplicity]\label{exr:iv12-expand-12}
What error arises if a double zero is counted only once in the argument principle?
\end{exercise}
\begin{hint}
The local residue equals the order.
\end{hint}
\begin{solution}
The count is off by one because a double zero contributes two, not one.
\end{solution}

\subsection*{Applications and investigations}

\begin{exercise}[Root counting without solving]\label{exr:iv12-expand-13}
Why is the argument principle useful for a high degree polynomial?
\end{exercise}
\begin{hint}
It converts root counting into a contour integral or argument change.
\end{hint}
\begin{solution}
One can determine the number of roots in a region without finding their individual values, provided the boundary contains no roots.
\end{solution}

\begin{exercise}[Nyquist style interpretation]\label{exr:iv12-expand-14}
Interpret the argument principle geometrically in terms of the image of the boundary under \(f\).
\end{exercise}
\begin{hint}
Track how \(f(\gamma)\) winds around zero.
\end{hint}
\begin{solution}
The net winding of the image about zero equals the number of enclosed zeros minus poles, counted with multiplicity.
\end{solution}

\subsection*{Challenge problems}

\begin{exercise}[Meromorphic divisor count]\label{exr:iv12-expand-15}
Show that a meromorphic function on a bounded region with no boundary zeros or poles has total divisor degree equal to the winding of its boundary values about zero.
\end{exercise}
\begin{hint}
Apply the argument principle.
\end{hint}
\begin{solution}
The divisor degree is the sum of zero orders minus pole orders. The argument principle identifies this integer with \((2\pi i)^{-1}\int f'/f\), which is the winding number of the image curve.
\end{solution}

\begin{exercise}[Nested contours]\label{exr:iv12-expand-16}
Two nested contours contain no zeros or poles in the annulus between them. Compare their argument-principle counts.
\end{exercise}
\begin{hint}
Apply the residue theorem to \(f'/f\) on the annulus.
\end{hint}
\begin{solution}
The two counts are equal because the difference of the contour integrals is the integral over the annular boundary, which vanishes when no zeros or poles lie between them.
\end{solution}

% END VOL04-EXPANSION IV12-exercises-01
